\documentclass{article}

% packages
\usepackage{amsmath, amsthm, thmtools, amsfonts, amssymb, luacode, catchfile, tikzducks, hyperref, ifthen}
\ifcsname c@kobocompile\endcsname
	\usepackage[a5paper, total={1072pt, 1448pt}, margin=10pt, includeheadfoot]{geometry} % set page margins
\else
	\usepackage[a4paper, margin=50pt, includeheadfoot]{geometry}
\fi
\usepackage[shortlabels]{enumitem}
\usepackage[skip=3pt, indent=0pt]{parskip}

% language
\usepackage[bidi=basic, layout=tabular, provide=*]{babel}
\ifcsname c@english\endcsname
	\babelprovide[main, import]{english}
\else
	\babelprovide[main, import]{hebrew}
	\babelprovide{rl}
\fi
%\babelfont{rm}{Libertinus Serif}
\babelfont{rm}[Renderer=Harfbuzz]{Libertinus Serif}
\babelfont{sf}{Libertinus Sans}
\babelfont{tt}{Libertinus Mono}

% style
\AddToHook{cmd/section/before}{\clearpage}	% Add line break before section
\linespread{1.3}
\setcounter{secnumdepth}{0}		% Remove default number tags from sections, this won't do well with theorems
\AtBeginDocument{\setlength{\belowdisplayskip}{3pt}}
\AtBeginDocument{\setlength{\abovedisplayskip}{3pt}}
\graphicspath{ {../images/} }

% operators
\DeclareMathOperator\cis{cis}
\DeclareMathOperator\Sp{Sp}
\DeclareMathOperator\tr{tr}
\DeclareMathOperator\im{Im}
\DeclareMathOperator\re{Re}
\DeclareMathOperator\diag{diag}
\DeclareMathOperator*\lowlim{\underline{lim}}
\DeclareMathOperator*\uplim{\overline{lim}}
\DeclareMathOperator\rng{rng}
\DeclareMathOperator\Sym{Sym}
\DeclareMathOperator\Arg{Arg}
\DeclareMathOperator\Log{Log}
\DeclareMathOperator\dom{dom}
\DeclareMathOperator\supp{Supp}
\DeclareMathOperator\var{Var}
\DeclareMathOperator\cov{Cov}

% commands
%\renewcommand\qedsymbol{\textbf{מש''ל}}
%\renewcommand\qedsymbol{\fbox{\emoji{lizard}}}
\newcommand{\Aa}[0]{\mathcal{A}}
\newcommand{\Bb}[0]{\mathcal{B}}
\newcommand{\CC}[0]{\mathbb{C}}
\newcommand{\Cc}[0]{\mathcal{C}}
\newcommand{\EE}[0]{\mathbb{E}}
\newcommand{\FF}[0]{\mathbb{F}}
\newcommand{\Ff}[0]{\mathcal{F}}
\newcommand{\Ii}[0]{\mathcal{I}}
\newcommand{\Gg}[0]{\mathcal{G}}
\newcommand{\Ll}[0]{\mathcal{L}}
\newcommand{\Mm}[0]{\mathcal{M}}
\newcommand{\NN}[0]{\mathbb{N}}
\newcommand{\Nn}[0]{\mathcal{N}}
\newcommand{\PP}[0]{\mathbb{P}}
\newcommand{\Pp}[0]{\mathcal{P}}
\newcommand{\QQ}[0]{\mathbb{Q}}
\newcommand{\RR}[0]{\mathbb{R}}
\newcommand{\Rr}[0]{\mathcal{R}}
\newcommand{\Ss}[0]{\mathcal{S}}
\newcommand{\TT}[0]{\mathbb{T}}
\newcommand{\Uu}[0]{\mathcal{U}}
\newcommand{\Vv}[0]{\mathcal{V}}
\newcommand{\Ww}[0]{\mathcal{W}}
\newcommand{\ZZ}[0]{\mathbb{Z}}
\newcommand{\acts}[0]{\circlearrowright}
\newcommand{\explain}[2] {
	\begin{flalign*}
		 && \text{#2} && \text{#1}
	\end{flalign*}
}
\newcommand{\maketitleprint}[0]{ \begin{center}
	%\begin{tikzpicture}[scale=3]
	%	\duck[graduate=gray!20!black, tassel=red!70!black]
	%\end{tikzpicture}	
	\includegraphics[width=6cm]{cover}
\end{center}
}

% theorem commands
\newtheoremstyle{c_remark}
	{}	% Space above
	{}	% Space below
	{}% Body font
	{}	% Indent amount
	{\bfseries}	% Theorem head font
	{}	% Punctuation after theorem head
	{.5em}	% Space after theorem head
	{\thmname{#1}\thmnumber{ #2}\thmnote{ \normalfont{\text{(#3)}}}}	% head content
\newtheoremstyle{c_definition}
	{3pt}	% Space above
	{3pt}	% Space below
	{}% Body font
	{}	% Indent amount
	{\bfseries}	% Theorem head font
	{}	% Punctuation after theorem head
	{.5em}	% Space after theorem head
	{\thmname{#1}\thmnumber{ #2}\thmnote{ \normalfont{\text{(#3)}}}}	% head content
\newtheoremstyle{c_plain}
	{3pt}	% Space above
	{3pt}	% Space below
	{\itshape}% Body font
	{}	% Indent amount
	{\bfseries}	% Theorem head font
	{}	% Punctuation after theorem head
	{.5em}	% Space after theorem head
	{\thmname{#1}\thmnumber{ #2}\thmnote{ \text{(#3)}}}	% head content

\ifcsname c@english\endcsname
	\theoremstyle{plain}
	\newtheorem{theorem}{Theorem}[section]
	\newtheorem{lemma}[theorem]{Lemma}
	\newtheorem{proposition}[theorem]{Proposition}
	\newtheorem*{proposition*}{Proposition}
	%\newtheorem{corollary}[theorem]{אין חלופה עברית}

	\theoremstyle{definition}
	\newtheorem{definition}[theorem]{Definition}
	\newtheorem*{definition*}{Definition}
	\newtheorem{example}{Example}[section]
	\newtheorem{exercise}{Exercise}[section]

	\theoremstyle{remark}
	\newtheorem*{remark}{Remark}
	\newtheorem*{solution}{Solution}
	\newtheorem{conclusion}[theorem]{Conclusion}
	\newtheorem{notation}[theorem]{Notation}
\else
	\theoremstyle{c_plain}
	\newtheorem{theorem}{משפט}[section]
	\newtheorem{lemma}[theorem]{למה}
	\newtheorem{proposition}[theorem]{טענה}
	\newtheorem*{proposition*}{טענה}
	%\newtheorem{corollary}[theorem]{אין חלופה עברית}

	\theoremstyle{c_definition}
	\newtheorem{definition}[theorem]{הגדרה}
	\newtheorem*{definition*}{הגדרה}
	\newtheorem{example}{דוגמה}[section]
	\newtheorem{exercise}{תרגיל}[section]

	\theoremstyle{c_remark}
	\newtheorem*{remark}{הערה}
	\newtheorem*{solution}{פתרון}
	\newtheorem{conclusion}[theorem]{מסקנה}
	\newtheorem{notation}[theorem]{סימון}
\fi

% Questions related commands
\newcounter{question}
\setcounter{question}{1}
\newcounter{sub_question}
\setcounter{sub_question}{1}

\ifcsname c@english\endcsname
	\newcommand{\question}[1][0]{
		\ifthenelse{#1 = 0}{}{\setcounter{question}{#1}}
		\section{Question \arabic{question}}
		\addtocounter{question}{1}
		\setcounter{sub_question}{1}
	}

	\newcommand{\subquestion}[1][0]{
		\ifthenelse{#1 = 0}{}{\setcounter{sub_question}{#1}}
		\subsection{Part \alph{sub_question}}
		\addtocounter{sub_question}{1}
	}
\else
	\newcommand{\question}[1][0]{
		\ifthenelse{#1 = 0}{}{\setcounter{question}{#1}}
		\section{שאלה \arabic{question}}
		\addtocounter{question}{1}
		\setcounter{sub_question}{1}
	}

	\newcommand{\subquestion}[1][0]{
		\ifthenelse{#1 = 0}{}{\setcounter{sub_question}{#1}}
		\subsection{סעיף \localecounter{letters.gershayim}{sub_question}}
		\addtocounter{sub_question}{1}
	}
\fi

% import lua and start of document
\directlua{common = require ('../common')}

\GetEnv{AUTHOR}

% headers
\author{\AUTHOR}
\date\today

\title{פתרון מטלה 06 --- אנליזה פונקציונלית, 80417}

\begin{document}
\maketitle
\maketitleprint[teal]

\question{}
יהי $1 \le p < \infty$ ונגדיר את המרחב $l^p$ להיות מרחב הסדרות הממשיות ${\{ x_n \}}_{n = 1}^\infty \subseteq \RR$ כך שמתקיים $\sum_{n = 1}^\infty {| x_n |}^p < \infty$.
נגדיר גם $\lVert x \rVert_p = {\left(\sum_{n = 1}^\infty {| x_n |}^p \right)}^{1/p}$, ראינו בקורסים קודמים כי זוהי נורמה על $l^p$.

\subquestion{}
נראה שאם $p \ne 2$ אז $\lVert \cdot \rVert_p$ אינה מושרית ממכפלה פנימית על $l^p$.
\begin{proof}
	נניח שקיימת מכפלה פנימית כך ש־$\sqrt{\langle x, x \rangle} = \lVert x \rVert_p$ לכל $x \in l^p$, אז נובע משאלה 4,
	\[
		2(\lVert x \rVert_p^2 + \lVert y \rVert_p^2)
		= \lVert x + y \rVert_p^2 + \lVert x - y \rVert_p^2
	\]
	ואילו נבחר $x_n = q^n$ אז,
	\[
		\lVert x \rVert_p^2
		= {\left(\sum_{n = 1}^\infty q^{np} \right)}^{2/p}
		= {\left( \frac{q^p}{1 - q^p} \right)}^{2/p}
		= \frac{q^2}{{(1 - q^p)}^{2/p}}
	\]
	בפרט עבור $q = \frac{1}{2}, \frac{1}{3}$,
	\[
		\lVert x + y \rVert_p^2
		= \frac{{(\frac{2}{3})}^2}{{(1 - {(\frac{2}{3})}^p)}^{2/p}},
		\qquad
		\lVert x - y \rVert_p^2
		= \frac{{(\frac{1}{6})}^2}{{(1 - {(\frac{1}{6})}^p)}^{2/p}}
	.\]
	ולכן,
	\begin{align*}
		& 2 \cdot \left(\frac{{(\frac{1}{3})}^2}{{(1 - {(\frac{1}{3})}^p)}^{2/p}} + \frac{{(\frac{1}{2})}^2}{{(1 - {(\frac{1}{2})}^p)}^{2/p}}\right)
		= \frac{{(\frac{2}{3})}^2}{{(1 - {(\frac{2}{3})}^p)}^{2/p}} + \frac{{(\frac{1}{6})}^2}{{(1 - {(\frac{1}{6})}^p)}^{2/p}} \\
		& \frac{9}{{(3^p - 1)}^{2/p}} + \frac{4}{{(2^p - 1)}^{2/p}} = \frac{18}{{(3^p - 2^p)}^{2/p}} + \frac{18}{{(6^p - 1)}^{2 / p}}
	\end{align*}
	ומבדיקה ישירה נקבל שרק $p = 2$ מקיים את הטענה, לכל ערך אחר מתקבלים באופן ישיר מספרים לא רציונליים ובפרט לא מתקבל פתרון של המשוואה.
	אבל הנחנו ש־$p \ne 2$ ולכן נסיק שאכן לא מתקיים חוק המקבילית, ואין מכפלה פנימית המשרה את הנורמה.

	נוכל לפתור בדרך נוספת.
	נגדיר $x = (1, 0, \ldots)$ וכן $y = (0, 1, 0, \ldots)$.
	במקרה זה $\lVert x \rVert_p^2 = \lVert y \rVert_p^2 = 1$ וכן $\lVert x + y \rVert_p^2 = \lVert x - y \rVert_p^2 = 2^{\frac{2}{p}}$.
	מכלל המקבילית נובע,
	\[
		2 (\lVert x \rVert_p^2 + \lVert y \rVert_p^2)
		= \lVert x + y \rVert_p^2 + \lVert x - y \rVert_p^2
		\iff 2 = 2^{\frac{2}{p}}
		\iff p = 2
	\]
	ונסיק שכלל המקבילית מתקיים אם ורק אם $p = 2$, לכן עבור $p \ne 2$ בהכרח הנורמה לא מושרית ממכפלה פנימית.
\end{proof}

\subquestion{}
נוכיח שנורמת $\infty$ על המרחב $C[a, b]$ לא מושרית ממכפלה פנימית.
\begin{proof}
	ללא הגבלת הכלליות נוכל לבחון את $C[0, 1]$, זאת שכן עבור ההוכחה שנציג אפשר להגדיר מתיחות של הפונקציות.
	נגדיר $f(x) = x, g(x) = 1 - x$ ונקבל $\lVert f \rVert_\infty^2 = \lVert g \rVert_\infty^2 = 1$.
	בנוסף גם $f + g = 1, f - g = 2x - 1$ ומתקיים $\lVert f + g \rVert_\infty^2 = \lVert f - g \rVert_\infty^2 = 1$.
	נסיק ישירות סתירה לכלל המקבילית.
\end{proof}

\question{}
יהי $(H, \langle \cdot, \cdot \rangle)$ מרחב מכפלה פנימית ויהיו $x \in H$ וכן ${\{ x_n \}}_{n = 1}^\infty \subseteq H$.

\subquestion{}
נראה כי אם $\lVert x_k \rVert \to \lVert x \rVert$ וגם $\langle x, x_k \rangle \to \langle x, x \rangle$ אז $x_k \to x$.
\begin{proof}
	מתקיים,
	\begin{align*}
		\lVert x - x_k \rVert
		& = \langle x - x_k, x - x_k \rangle \\
		& = \langle x - x_k, x \rangle + \langle x - x_k, -x_k \rangle \\
		& = \overline{\langle x, x - x_k \rangle} + \overline{\langle -x_k, x - x_k \rangle} \\
		& = \overline{\langle x, x \rangle} + \overline{\langle x, -x_k \rangle} - \overline{\langle x_k, x \rangle} - \overline{\langle x_k, -x_k \rangle} \\
		& = \lVert x \rVert^2 - \langle x_k, x \rangle - \overline{\langle x_k, x \rangle} + \lVert x_k \rVert^2 \\
		& \xrightarrow{k \to \infty} \lVert x \rVert^2 - \langle x, x \rangle - \langle x, x \rangle + \lVert x \rVert^2 \\
		& = 0
	\end{align*}
	ולכן נסיק ש־$x_k \to x$.
\end{proof}

\subquestion{}
תהי ${\{ e_k \}}_{k = 1}^\infty \subseteq H$ מערכת אורתונורמלית,
נראה שהסדרה $y_k = \langle e_k, x \rangle \cdot e_k$ מתכנסת.
\begin{proof}
	מאי־שוויון בסל נובע,
	\[
		\lVert x \rVert
		\ge \sum_{k = 1}^\infty {|\langle e_k, x \rangle|}^2
	\]
	ולכן מהתנאי ההכרחי להתכנסות בהחלט של טורים ממשיים נובע,
	\[
		{|\langle e_k, x \rangle|}^2
		\to 0
		\implies \langle e_k, x \rangle \to 0
	\]
	נסיק אם כן ש־$\lVert y_k \rVert = |\langle e_k, x \rangle| \to 0$.
	נבחין גם כי $\langle y_k, 0 \rangle = 0$ לכל $k \in \NN$ ישירות מהומוגניות מכפלה פנימית, ולכן תנאי סעיף א' חלים ו־$y_k \to 0$.
\end{proof}

\question{}
נוכיח כי הבסיס הסטנדרטי הוא מערכת אורתונורמלית שלמה ב־$l^2$.
\begin{proof}
	נסמן ${\{ e^k \}} \subseteq l^2$ הבסיס הסטנדרטי, כלומר $e_n^k = 1 \iff n = k$ ו־$e_n^k = 0$ אחרת. \\
	לכל $k \ne m$ מתקיים,
	\[
		\langle e^k, e^m \rangle
		= \sum_{n = 1}^\infty \overline{e_n^k} e_n^m
		= 0
	\]
	בנוסף מאותה סיבה בדיוק $\langle e^k, e^k \rangle = 1$ לכל $k \in \NN$.
	ונסיק כי זה אכן בסיס אורתונורמלי.

	נראה ש־$\overline{\Sp\{ e^k \mid k \in \NN \}} = l^2$ ונסיק משקילות שהבסיס שלם.
	נניח ש־$x \in l^2$ סדרה כלשהי, ונרצה לבנות סדרה ${\{ a^k \}} \subseteq \Sp\{ e^k \}$ כך ש־$a^k \to x$.
	נגדיר,
	\[
		a_n^k
		= \begin{cases}
			x_n & n \le k \\
			0 & n > k
		\end{cases}
	\]
	כלומר $a^k = \sum_{n = 1}^n x_n e^n$, לכן $a^k \in \Sp\{ e^k \}$ לכל $k \in \NN$.
	כמובן גם $a_n^k \to x$ שכן $\lVert x - a^k \rVert^2 = \sum_{n = k}^\infty x_n^k$.
	נסיק שאכן הבסיס אורתונורמלי ושלם.
\end{proof}

\question{}
יהי $(H, \langle \cdot, \cdot \rangle)$ מרחב מכפלה פנימית.

\subquestion{}
נראה כי נורמה המושרית ממכפלה פנימית מקיימת את חוק המקבילית, כלומר שלכל $x, y \in H$,
\[
	2( \lVert x \rVert^2 + \lVert y \rVert^2) = {\lVert x + y \rVert}^2 + {\lVert x - y \rVert}^2
\]
\begin{proof}
	יהיו $x, y \in H$, אז,
	\[
		2( \lVert x \rVert^2 + \lVert y \rVert^2) 
		= 2 \langle x, x \rangle + 2 \langle y, y \rangle
	\]
	וכן,
	\begin{align*}
		{\lVert x + y \rVert}^2 + {\lVert x - y \rVert}^2
		& = \langle x + y, x + y \rangle - \langle x - y, x - y \rangle \\
		& = \langle x + y, x \rangle + \langle x + y, y \rangle + \langle x - y, x \rangle - \langle x - y, y \rangle \\
		& = \overline{\langle x, x + y \rangle} + \overline{\langle y, x + y \rangle} + \overline{\langle x, x - y \rangle} - \overline{\langle y, x - y \rangle} \\
		& = \overline{\langle x, x \rangle} + \overline{\langle y, x \rangle} + \overline{\langle x, x \rangle} - \overline{\langle y, x \rangle}
		+ \overline{\langle x, y \rangle} + \overline{\langle y, y \rangle} - \overline{\langle x, y \rangle} + \overline{\langle y, y \rangle} \\
		& = 2 \langle x, x \rangle + 2 \langle y, y \rangle
	\end{align*}
	כפי שרצינו להראות.
\end{proof}

\subquestion{}
נניח ש־$H$ מרחב מכפלה פנימית ממשי, אנו נראה כי לכל $x, y \in H$,
\[
	\langle x, y \rangle
	= \frac{\lVert x + y \rVert^2 - \lVert x - y \rVert^2}{4}
\]
\begin{proof}
	זהו מרחב ממשי ולכן $\langle x, y \rangle = \langle y, x \rangle$, לכן,
	\[
		\lVert x + y \rVert^2 - \lVert x - y \rVert^2
		= {\lVert x \rVert}^2 + {\lVert y \rVert}^2 + 2 \langle x, y \rangle
		- ({\lVert x \rVert}^2 + {\lVert y \rVert}^2 - 2 \langle x, y \rangle)
		= 4 \cdot \langle x, y \rangle
	\]
	והנוסחה נובעת ישירות.
\end{proof}

\end{document}
